\label{ch1:introduction}

Chiplet-based design has become an important approach for building scalable and configurable hardware systems. Instead of implementing a complete system as one large monolithic chip, the design can be divided into several smaller chips and integrated in the same package. This approach can improve design reuse, yield, and flexibility, especially when different functions or processing elements need to be combined for a specific application.

However, as more components are integrated, data movement between them becomes an important design issue. Processing elements, memories, and accelerators need to exchange data efficiently, and the communication structure can strongly affect the overall system behavior. A simple shared bus may become inefficient when the number of connected units increases, while a fully connected structure can be too expensive in terms of hardware resources. Therefore, a scalable communication fabric is needed.

Network-on-Chip (NoC) is commonly used as a structured communication method for connecting multiple components in a chip or chiplet-based system. A NoC usually consists of routers, links, and network interfaces. Packets are forwarded through routers according to a routing algorithm, and the network can be extended by connecting more routers together. Compared with bus-based communication, NoC provides better scalability and can support multiple simultaneous data transfers. NoC is also widely considered a promising solution for addressing communication bottlenecks in large-scale multicore and heterogeneous systems \cite{10819599}.

\section{Related work}
Network-on-Chip has been widely studied as a communication fabric for manycore, heterogeneous, and chiplet-based systems. In heterogeneous systems, different types of processing elements, such as CPUs, GPUs, memories, and accelerators, may have different communication requirements. For example, CPU-related traffic is often latency-sensitive, while GPU-related traffic usually requires higher throughput. This makes the design of the on-chip communication network more challenging than in a homogeneous system. Hybrid NoC architectures combining wireline and wireless links have been explored to improve communication efficiency in CPU-GPU heterogeneous manycore systems, especially for workloads with intensive data movement such as CNN training \cite{8119941}.

NoC design has also been investigated in the context of chiplet-based manycore architectures. Flexible NoC architectures with adaptable routers and links have been proposed to dynamically form different sub-network topologies, such as mesh, concentrated mesh, torus, and tree \cite{9218654}. These topologies can be selected according to application traffic patterns and different latency or bandwidth requirements. Such studies show that adapting the network structure can improve system performance and energy efficiency. However, this flexibility also introduces additional hardware and control complexity.

\section{Thesis outline}
The goal of this thesis is to design and implement a NoC router for efficient data movement in a small mesh-based communication system. The router is intended to support packet-based communication between processing elements and to provide basic routing, virtual-channel flow control, and pipelined forwarding. A 3×3 mesh network is used as the target structure, and FPGA implementation is considered as the hardware platform for evaluating the design.

The thesis work is organized as follows:

\begin{itemize}
    \item \textbf{Chapter 2: Background.} 
    Introduces the key concepts used in this work, including chiplet-based systems, 
    Network-on-Chip, mesh topology, packet switching, routing, virtual channels, 
    and flow control.

    \item \textbf{Chapter 3: Design and Implementation.} 
    Describes the proposed router architecture and its main design components. 
    This chapter explains the flit format, routing algorithm, allocation methods, 
    pipeline organization, and 3$\times$3 mesh integration.

    \item \textbf{Chapter 4: Verification and Results.} 
    Presents the verification plan and FPGA-oriented implementation method. 
    It describes single-router verification, mesh-level verification, test cases, 
    and the use of UART for test control and result reporting.

    \item \textbf{Chapter 5: Conclusion and Future Work.} 
    Summarizes the thesis work and discusses possible future improvements, such as 
    extending the design to different network topologies, exploring adaptive routing 
    algorithms, improving arbitration schemes, and integrating the NoC with real 
    processing elements or accelerator modules.
\end{itemize}


   
